\documentclass{article}

\usepackage[margin=1in]{geometry}
\usepackage{amsmath,amsthm,amssymb}
\usepackage[spanish]{babel}
\usepackage{enumitem}

\theoremstyle{plain}
\newtheorem{proposicion}{Proposición}[section]

\theoremstyle{definition}
\newtheorem{definition}{Definición}[section]

% Number sets
\newcommand{\R}{\mathbb{R}}
\newcommand{\Z}{\mathbb{Z}}
\newcommand{\N}{\mathbb{N}}
\newcommand{\Q}{\mathbb{Q}}
\newcommand{\C}{\mathbb{C}}

% Topology operators
\newcommand{\cl}[1]{\overline{#1}}              % Closure
\newcommand{\Int}[1]{{#1}^\circ}                % Interior
\newcommand{\bd}{\partial}                      % Boundary
\newcommand{\Ent}[1]{\mathcal{O}(#1)}           % Neighborhoods (Entornos)
\newcommand{\pf}{\mathcal{P}_F}                 % Finite parts
\newcommand{\partes}{\mathcal{P}}               % Power set

% Useful shortcuts
\newcommand{\imp}{\implies}                      % Implies

% Net convergence/accumulation notation
\newcommand{\evto}[1]{\stackrel{E}{\hookrightarrow} #1}   % Eventually in
\newcommand{\frto}[1]{\stackrel{F}{\hookrightarrow} #1}   % Frequently in
\newcommand{\nnevto}[1]{\not\!\stackrel{E}{\hookrightarrow} #1}
\newcommand{\nnfrto}[1]{\not\!\stackrel{F}{\hookrightarrow} #1}

% Exercise environment
\newenvironment{ejercicio}[1]{\begin{trivlist}
\item[\hskip \labelsep {\bfseries Ejercicio}\hskip \labelsep {\bfseries #1.}]}{\end{trivlist}}

\setlist[enumerate,1]{label=(\alph*), leftmargin=*, itemsep=2pt}

\begin{document}

\title{Topología: Práctica VII}
\author{Bustos Jordi\\Espacios localmente compactos}

\maketitle

\begin{ejercicio}{1}
    Sea \( X \) un espacio topológico LKH. Probar que:
    \begin{enumerate}
        \item[(a)] \( C_{K}(X,\mathbb{C}) \subseteq C_{0}(X, \mathbb{C}) \subseteq C_{b}(X, \mathbb{C}) \) y son métricos con la \( d_{\infty} \).
        \item[(b)] \( C_{K}(X,\mathbb{C}) \) es \( d_{\infty} \)-denso en \( C_{0}(X,\mathbb{C}) \).
        \item[(c)] Dada \( f \in C_{0}(X,\mathbb{C}) \) y \( \varepsilon > 0 \), el conjunto \( \{x \in X : |f(x)| \ge \varepsilon \} \) es compacto en \( X \).
        \item[(d)] Si \( X \) es métrico entonces toda \( f \in C_{0}(X,\mathbb{C}) \) es uniformemente continua.
    \end{enumerate}
\end{ejercicio}

\begin{proof}
    Para (a), dada \( f \in C_{K}(X, \mathbb{C}) \) existe un compacto \( K \subseteq X \) tal que \( f(x) = 0 \) para todo \( x \in X \setminus K \). Por lo tanto, para todo \( \varepsilon > 0 \) si tomo \( K_\varepsilon = K \), entonces
    \( |f(x)| = 0 < \varepsilon \) para todo \( x \notin K_\varepsilon \), es decir, \( f \in C_{0}(X, \mathbb{C}) \).

    Además, \( C_{0}(X, \mathbb{C}) \subseteq C_{b}(X, \mathbb{C}) \) porque si \( f \in C_{0}(X, \mathbb{C}) \) entonces existe un compacto \( K_1 \subseteq X \) tal que \( |f(x)| < 1 \) para todo \( x \notin K_1 \).
    Como \( f \) es continua y \( K_1 \) es compacto, alcanza un máximo en \( K_1 \). Por lo tanto, \( |f(x)| < 1 + M \) para todo \( x \in X \), donde \( M = \max_{x\in K_1} |f(x)| < +\infty \).

    Finalmente, la métrica \( d_{\infty} \) es la misma en los tres espacios y los hace métricos. En efecto, sean \( f, g \in C_b(X, \mathbb{C}) \), \( d_{\infty}(f,g) = \sup_{x \in X} |f(x) - g(x)| \), luego:
    \[
        d_{\infty}(f,g) = 0 \iff f(x) = g(x) \quad \forall x \in X \iff f = g
    \]
    Además, \( d_{\infty}(f,g) = d_{\infty}(g,f) \) y por desigualdad triangular en \( \mathbb{C} \), \( d_{\infty}(f,h) \le d_{\infty}(f,g) + d_{\infty}(g,h) \) para todo \( f, g, h \in C_b(X, \mathbb{C}) \).

    Para el inciso (b), dado \( f \in C_0(X, \mathbb{C}) \) y \( \varepsilon > 0 \), existe un compacto \( K_\varepsilon \subseteq X \) tal que \( |f(x)| < \varepsilon \) para todo \( x \notin K_\varepsilon \).
    Como \( X \) es un espacio LKH, por el Lema de Urysohn existe una función continua con soporte compacto \( h: X \to [0,1] \) tal que \( h(x) = 1 \) para todo \( x \in K_\varepsilon \).
    Definimos \( g = f \cdot h \). Como \( g \) es producto de continuas es continua, y al ser \( h \) de soporte compacto, \( g \) también lo es, luego \( g \in C_K(X, \mathbb{C}) \). Evaluemos la distancia:
    \[
        d_{\infty}(f,g) = \sup_{x\in X} |f(x) - f(x)h(x)| = \sup_{x\in X} |f(x)| |1 - h(x)|
    \]
    Si \( x \in K_\varepsilon \), \( h(x) = 1 \implies |f(x)||1 - h(x)| = 0 \). Si \( x \notin K_\varepsilon \), sabemos que \( |f(x)| < \varepsilon \) y como \( 0 \le h(x) \le 1 \), entonces \( |1 - h(x)| \le 1 \). Por lo tanto:
    \[
        d_{\infty}(f,g) \le \sup_{x \notin K_\varepsilon} |f(x)| < \varepsilon
    \]
    Concluyendo que \( C_K(X, \mathbb{C}) \) es \( d_{\infty} \)-denso en \( C_0(X, \mathbb{C}) \).

    Para el inciso (c), dada \( f \in C_0(X, \mathbb{C}) \) y \( \varepsilon > 0 \), existe un compacto \( K_\varepsilon \subseteq X \) tal que \( |f(x)| < \varepsilon \) para todo \( x \notin K_\varepsilon \).
    Esto implica que el conjunto \( K = \{ x \in X : |f(x)| \geq \varepsilon \} \) está contenido en \( K_\varepsilon \). Como \( f \) es continua, \( K \) es la preimagen del conjunto cerrado \( \{ z \in \mathbb{C} : |z| \ge \varepsilon \} \), por lo que \( K \) es cerrado en \( X \). Al ser un subconjunto cerrado dentro del compacto \( K_\varepsilon \), \( K \) es compacto.

    Para el inciso (d), sea \( X \) métrico y \( f \in C_0(X, \mathbb{C}) \). Dado \( \varepsilon > 0 \), por el inciso anterior sabemos que el conjunto \( K = \{ x \in X : |f(x)| \ge \varepsilon/2 \} \) es compacto.
    Por la continuidad de \( f \), para cada \( x \in K \) existe un radio \( \delta_x > 0 \) tal que si \( d(y,x) < 2\delta_x \), entonces \( |f(y) - f(x)| < \varepsilon/2 \).
    La familia de bolas \( {\{ B(x, \delta_x) \}}_{x \in K} \) es un cubrimiento abierto del compacto \( K \), por lo que admite un subcubrimiento finito \( B(x_1, \delta_1), \dots, B(x_n, \delta_n) \).
    Tomemos \( \delta = \min \{ \delta_1, \dots, \delta_n \} > 0 \). Sean \( y, z \in X \) con \( d(y,z) < \delta \). Tenemos dos casos:
    \begin{itemize}
        \item Caso 1: Ninguno de los puntos pertenece a \( K \). Entonces \( |f(y)| < \varepsilon/2 \) y \( |f(z)| < \varepsilon/2 \), lo que implica \( |f(y) - f(z)| \le |f(y)| + |f(z)| < \varepsilon \).
        \item Caso 2: Al menos uno de los puntos pertenece a \( K \). Supongamos sin pérdida de generalidad que \( y \in K \). Entonces existe algún \( i \in \{1,\dots,n\} \) tal que \( y \in B(x_i, \delta_i) \), es decir \( d(y, x_i) < \delta_i \).
              Por la desigualdad triangular, \( d(z, x_i) \le d(z, y) + d(y, x_i) < \delta + \delta_i \le 2\delta_i \).
              Como tanto \( y \) como \( z \) están a distancia menor a \( 2\delta_i \) de \( x_i \), evaluando en la función obtenemos \( |f(y) - f(x_i)| < \varepsilon/2 \) y \( |f(z) - f(x_i)| < \varepsilon/2 \). Sumando ambas desigualdades resulta \( |f(y) - f(z)| < \varepsilon \).
    \end{itemize}
    En cualquier caso, si \( d(y,z) < \delta \), entonces \( |f(y) - f(z)| < \varepsilon \), lo que demuestra que \( f \) es uniformemente continua en \( X \).
\end{proof}

\begin{ejercicio}{2}
    Sean \( X \) e \( Y \) dos espacios topológicos. Considerar en \( C(X,Y) \) la topología compacto-abierta. Demostrar que si \( Y \) es Hausdorff (o regular) entonces \( C(X,Y) \) es Hausdorff (o regular).
\end{ejercicio}
\begin{proof}

\end{proof}

\begin{ejercicio}{3}
    Sea \( (X, \tau) \) un espacio topológico. Probar que si \( X \) es compacto entonces es PK. Además, si \( X \) es unión disjunta de subconjuntos abiertos y PK's, entonces \( X \) es PK.
\end{ejercicio}
\begin{proof}

\end{proof}

\begin{ejercicio}{4}
    Sea \( (X, \tau) \) espacio PK. Probar lo siguiente:
    \begin{enumerate}
        \item[(a)] Todo \( F \subseteq X \) que sea cerrado queda PK con la topología inducida.
        \item[(b)] Si \( Y \) es otro espacio que es compacto entonces \( X \times Y \) es PK.
        \item[(c)] En general, \( X \times X \) no es PK.
    \end{enumerate}
\end{ejercicio}
\begin{proof}

\end{proof}

\begin{ejercicio}{5}
    Sea \( X \) un espacio topológico LKH. Dada una familia \( \mathcal{F} \subseteq C_{0}(X) \) se tiene que \( \mathcal{F} \) es totalmente acotada \( \iff \) \( \mathcal{F} \) es acotada, equicontinua y se anula uniformemente en el \( \infty \).
\end{ejercicio}
\begin{proof}

\end{proof}

\begin{ejercicio}{6}
    Sea \( X \) un espacio topológico LK y \( f \in C(X,Y) \) suryectiva y abierta. Probar que:
    \begin{enumerate}
        \item[(a)] \( Y \) es LK.
        \item[(b)] Si \( Y \) es Hausdorff entonces todo compacto \( C \) de \( Y \) es de la forma \( f(K)=C \) para algún \( K \) compacto de \( X \).
    \end{enumerate}
\end{ejercicio}
\begin{proof}

\end{proof}

\begin{ejercicio}{7}
    Sea \( X \) un PKH. Sea \( \mathcal{C} \) una familia de conjuntos de \( X \) y para cada \( C \in \mathcal{C} \) sea un número \( \varepsilon_{C} > 0 \). Si \( \mathcal{C} \) es LF entonces existe \( f:X \rightarrow \mathbb{R} \) continua tal que \( f(x) > 0 \) para todo \( x \in X \) y \( f(x) \le \varepsilon_{C} \) para todo \( x \in C \).
\end{ejercicio}
\begin{proof}

\end{proof}

\begin{ejercicio}{8}
    Dado \( A \) no numerable, mostrar que no es PK, considerando el producto de las topologías discretas \( \prod_{a \in A} \mathbb{N} \).
\end{ejercicio}
\begin{proof}

\end{proof}

\end{document}